\subsection{问题二：变物性热质耦合模型的建立与求解}
\subsubsection{问题二模型的建立}
问题二从$t=0$起统一使用题目附录3的经验式：
\begin{equation}
\begin{aligned}
\rho(C)&=650+128C,\\
c_p(C)&=1450+2736\frac{C}{1+C},\\
k(C)&=0.21+0.38\frac{C}{1+C},\\
D(C,T_K)&=2.4\times10^{-3}
\exp\left(-\frac{0.45}{C}-\frac{3850}{T_K}\right).
\end{aligned}
\label{eq:props2}
\end{equation}

密度和比热容共同控制单位体积升温所需热量，导热系数影响内部热流，扩散系数影响水分迁移。将其代入守恒形式得到
\begin{equation}
\begin{cases}
\rho(C)c_p(C)T_t=\dfrac1r\pd_r[rk(C)T_r],\\[4pt]
C_t=\dfrac1r\pd_r[rD(C,T_K)C_r].
\end{cases}
\label{eq:q2pde}
\end{equation}

初始状态仍为$T=28$、$C=2.55$；中心对称条件和表面对流形式保持不变，表面导热与扩散系数按当前表面状态取值。这样便形成物性随解更新的非线性初边值问题。

对式\eqref{eq:props2}求偏导可得
\begin{equation}
D_C=\frac{0.45}{C^2}D>0,\qquad
D_T=\frac{3850}{T_K^2}D>0.
\end{equation}

温度升高会增强扩散，含水率降低则减弱扩散。空间展开进一步得到
\begin{equation}
\begin{aligned}
\frac1r\pd_r(rDC_r)
&=D\left(C_{rr}+\frac{C_r}{r}\right)
+D_C(C_r)^2+D_TT_rC_r,\\
\frac1r\pd_r(rkT_r)
&=k\left(T_{rr}+\frac{T_r}{r}\right)+k_C C_rT_r .
\end{aligned}
\end{equation}


\subsubsection{问题二模型的求解}
问题二继续采用式\eqref{eq:mesh}的表面加密网格和式\eqref{eq:ode}的控制体更新。区别在于，每次右端评估都要重新计算节点热容量及界面$k,D$。在隐式迭代过程中，温度与含水率同时变化，最终接受的时间步应满足同一组耦合残差方程，避免使用相互不一致的旧物性。

主计算$N=2560,\ \beta=5$，相对容差$2\times10^{-10}$，前4 h最大步长30 s，后期上限900 s；绝对容差按温度和含水率各自尺度设置。积分器在变化剧烈或数据折点附近自动缩短步长。后期上限较大并不代表实际每900 s仅计算一次，因为步长仍受误差估计约束。

\fig{q2_summary}{0.98}{由问题二已保存结果提取的代表位置演化；标记为0.5 h间隔表值}{fig:q2summary}

图\ref{fig:q2summary}反映两种不同的时间尺度。温度在前3 h逐渐靠近烘房稳定温度，中心和表面的差距持续缩小；含水率虽整体下降，但3 h时仍远高于干燥要求。较快的热响应不能直接转化为同样快的干燥完成，因为水分迁移具有独立的扩散阻力，并随含水率降低而减慢。


\begin{table}[H]\centering
\caption{问题2：温度分布（$^\circ$C）}\label{tab:q2T}
\begin{tabular}{cccccc}\toprule
时间/h & 0 cm & 0.5 cm & 1 cm & 1.5 cm & 2 cm\\ \midrule
0.5 & 32.1892 & 32.3821 & 32.9659 & 33.9608 & 35.4130\\
1.0 & 40.3816 & 40.5536 & 41.0605 & 41.8782 & 42.9976\\
1.5 & 45.8468 & 45.9348 & 46.1930 & 46.6051 & 47.1401\\
2.0 & 48.4502 & 48.4881 & 48.5984 & 48.7736 & 49.0033\\
2.5 & 49.4670 & 49.4792 & 49.5136 & 49.5653 & 49.6609\\
3.0 & 49.8495 & 49.8553 & 49.8746 & 49.9101 & 49.9664\\
\bottomrule\end{tabular}\end{table}

\begin{table}[H]\centering
\caption{问题2：干基含水率分布（kg/kg）}\label{tab:q2C}
\begin{tabular}{cccccc}\toprule
时间/h & 0 cm & 0.5 cm & 1 cm & 1.5 cm & 2 cm\\ \midrule
0.5 & 2.5499 & 2.5489 & 2.5256 & 2.3257 & 1.6486\\
1.0 & 2.5257 & 2.4948 & 2.3578 & 2.0230 & 1.4711\\
1.5 & 2.3861 & 2.3256 & 2.1344 & 1.8020 & 1.3476\\
2.0 & 2.1709 & 2.1084 & 1.9236 & 1.6259 & 1.2311\\
2.5 & 1.9566 & 1.9006 & 1.7360 & 1.4720 & 1.1166\\
3.0 & 1.7662 & 1.7165 & 1.5702 & 1.3333 & 1.0081\\
\bottomrule\end{tabular}\end{table}


由表\ref{tab:q2T}可见，0.5 h时中心温度32.1892 $^\circ$C、表面35.4130 $^\circ$C；3 h时分别为49.8495 $^\circ$C和49.9664 $^\circ$C，中心与表面温差已缩小到0.1169 $^\circ$C。此时热量在径向上的再分配基本完成，但仍保持小幅传热驱动力。

表\ref{tab:q2C}显示，3 h时中心和表面含水率分别为1.7662、1.0081 kg/kg，约为目标阈值的11.8倍和6.7倍。故温度接近环境并不是干燥结束标准。若仅根据温度平稳判断结束，会明显低估后期失水所需时间。

问题一与问题二在同样0.5 h时刻得到不同值，原因是两问从开始就采用不同物性与扩散系数，而不是单位换算或初值衔接错误。比较两问应理解为不同参数模型的响应比较，不应把两张表首尾连接成实际阶段切换轨迹。若要模拟真实的分阶段切换，还需给出切换时刻及参数过渡规则，现有题目要求则是统一使用附录3。